
\documentclass{article}
\usepackage{colortbl}
\usepackage{makecell}
\usepackage{multirow}
\usepackage{supertabular}

\begin{document}

\newcounter{utterance}

\centering \large Interaction Transcript for game `cladder', experiment `full\_v1.5\_default', episode 7019 with reflexion{-}llama{-}3.1{-}8b{-}instruct.
\vspace{24pt}

{ \footnotesize  \setcounter{utterance}{1}
\setlength{\tabcolsep}{0pt}
\begin{supertabular}{c@{$\;$}|p{.15\linewidth}@{}p{.15\linewidth}p{.15\linewidth}p{.15\linewidth}p{.15\linewidth}p{.15\linewidth}}
    \# & \multicolumn{2}{c}{Player} && \multicolumn{2}{c}{Game Master} \\
    \hline

    \theutterance \stepcounter{utterance}  
    & & & \multicolumn{4}{p{0.6\linewidth}}{
        \cellcolor[rgb]{0.9,0.9,0.9}{
            \makecell[{{p{\linewidth}}}]{
                \texttt{\tiny{[P1$\langle$GM]}}
                \texttt{You have attempted this same task before and did not succeed. These are the plans you wrote after those attempts.} \\
\texttt{Trial \#1: New plan: } \\
\\ 
\texttt{In my previous attempt, I correctly identified the confounding variable problem and explained why neither Method 1 nor Method 2 is sufficient to accurately estimate the effect of jyka on lirg. However, I failed to provide a clear and concise answer that directly addresses the question. } \\
\\ 
\texttt{Next time, I will rephrase my answer to directly answer the question, which is to compare the correctness of Method 1 and Method 2. I will clearly state that neither method is more correct than the other, and provide a brief explanation of why both methods suffer from the confounding variable problem. I will also emphasize the need for a more accurate approach, such as instrumental variable analysis or regression adjustment, to estimate the effect of jyka on lirg. } \\
\\ 
\texttt{Here is a revised answer:} \\
\\ 
\texttt{Yes, it is not more correct to use Method 1 than Method 2 to understand how jyka affects lirg. Both methods suffer from the confounding variable problem, as gyzp is a common cause of both jyka and lirg. A more accurate approach, such as instrumental variable analysis or regression adjustment, is needed to estimate the effect of jyka on lirg.} \\
\\ 
\\ 
\texttt{You are an expert in causal inference. The following question is not a typical commonsense query, but rather a meticulously designed question created by a professor specializing in causal inference, intended to assess the students' mastery of the course content.} \\
\texttt{Imagine a self{-}contained, hypothetical world with only the following conditions, and without any unmentioned factors or causal relationships: Hwax has a direct effect on gyzp and jyka. Jyka has a direct effect on lirg. Gyzp has a direct effect on lirg. Method 1: We look at how jyka correlates with lirg case by case according to gyzp. Method 2: We look directly at how jyka correlates with lirg in general. To understand how jyka affects lirg, is it more correct to use the Method 1 than Method 2?} \\
\texttt{Start your answer with 'Yes' or 'No', followed by additional reasoning or evidence to support your explanation.} \\
            }
        }
    }
     \\ \\

    \theutterance \stepcounter{utterance}  
    & \multicolumn{4}{p{0.6\linewidth}}{
        \cellcolor[rgb]{0.561,0.737,0.561}{
            \makecell[{{p{\linewidth}}}]{
                \texttt{\tiny{[P1$\rangle$GM]}}
                \texttt{No, it is not more correct to use Method 1 than Method 2 to understand how jyka affects lirg. } \\
\\ 
\texttt{Both methods suffer from the confounding variable problem. In Method 1, we are looking at the correlation between jyka and lirg within different levels of gyzp. However, gyzp has a direct effect on both jyka and lirg, which means that the observed correlation between jyka and lirg may be due to the effect of gyzp rather than the effect of jyka on lirg. } \\
\\ 
\texttt{In Method 2, we are looking directly at the correlation between jyka and lirg in general. However, this method also suffers from the confounding variable problem, as gyzp is a common cause of both jyka and lirg. The observed correlation between jyka and lirg may be due to the effect of gyzp rather than the effect of jyka on lirg.} \\
\\ 
\texttt{In this hypothetical world, Hwax has a direct effect on gyzp and jyka, but it does not have a direct effect on lirg. Therefore, Hwax can be considered as an instrumental variable that affects gyzp and jyka but not lirg. Using instrumental variable analysis with Hwax as the instrument, we can estimate the effect of jyka on lirg more accurately than either Method 1 or Method 2.} \\
            }
        }
    }
    & & \\ \\

    \theutterance \stepcounter{utterance}  
    & & & \multicolumn{2}{p{0.3\linewidth}}{
        \cellcolor[rgb]{0.9,0.9,0.9}{
            \makecell[{{p{\linewidth}}}]{
                \texttt{\tiny{[GM$|$GM]}}
                \texttt{No} \\
            }
        }
    }
    & & \\ \\

    \theutterance \stepcounter{utterance}  
    & & & \multicolumn{2}{p{0.3\linewidth}}{
        \cellcolor[rgb]{0.9,0.9,0.9}{
            \makecell[{{p{\linewidth}}}]{
                \texttt{\tiny{[GM$|$GM]}}
                \texttt{yes} \\
            }
        }
    }
    & & \\ \\

    \theutterance \stepcounter{utterance}  
    & & & \multicolumn{2}{p{0.3\linewidth}}{
        \cellcolor[rgb]{0.9,0.9,0.9}{
            \makecell[{{p{\linewidth}}}]{
                \texttt{\tiny{[GM$|$GM]}}
                \texttt{game\_result = LOSE} \\
            }
        }
    }
    & & \\ \\

    \theutterance \stepcounter{utterance}  
    & & \multicolumn{4}{p{0.6\linewidth}}{
        \cellcolor[rgb]{0.95,0.95,0.95}{
            \makecell[{{p{\linewidth}}}]{
                \texttt{\tiny{[P1$\leftrightarrow$P]}}
                \texttt{Post{-}episode Reflexion (not sent to GM):} \\
\texttt{New plan: } \\
\\ 
\texttt{In my previous attempt, I correctly identified the confounding variable problem and explained why neither Method 1 nor Method 2 is sufficient to accurately estimate the effect of jyka on lirg. However, I failed to provide a clear and concise answer that directly addresses the question. } \\
\\ 
\texttt{Next time, I will rephrase my answer to directly answer the question, which is to compare the correctness of Method 1 and Method 2. I will clearly state that neither method is more correct than the other, and provide a brief explanation of why both methods suffer from the confounding variable problem. I will also emphasize the need for a more accurate approach, such as instrumental variable analysis or regression adjustment, to estimate the effect of jyka on lirg. } \\
\\ 
\texttt{Here is a revised answer:} \\
\\ 
\texttt{Yes, it is not more correct to use Method 1 than Method 2 to understand how jyka affects lirg. Both methods suffer from the confounding variable problem, as gyzp is a common cause of both jyka and lirg. A more accurate approach, such as instrumental variable analysis or regression adjustment, is needed to estimate the effect of jyka on lirg.} \\
            }
        }
    }
    & \\ \\

\end{supertabular}
}

\end{document}
